\documentclass[12pt, a4paper]{article}
\usepackage{graphicx} % Required for inserting images
\graphicspath{{images/}{images-desafios/}}
\usepackage[a4paper, total={7.5in, 10in}]{geometry}
\usepackage{float} % pra conseguir usar o \begin{figure}[H]
\usepackage{indentfirst} % pra fazer tab n 1 paragrafos dos capitulos
% package pra tabela
\usepackage{multirow}
\usepackage{array}

%pacote pra codigo
\usepackage{listings} 
\usepackage{xcolor}
\definecolor{codegreen}{rgb}{0,0.6,0}
\definecolor{codegray}{rgb}{0.5,0.5,0.5}
\definecolor{codepurple}{rgb}{0.58,0,0.82}
\definecolor{backcolour}{rgb}{0.95,0.95,0.92}

\lstdefinestyle{myCodeStyle}{
	backgroundcolor=\color{backcolour},   
	commentstyle=\color{codegreen},
	keywordstyle=\color{cyan},
	numberstyle=\tiny\color{codegray},
	stringstyle=\color{codepurple},
	basicstyle=\ttfamily\footnotesize,
	breakatwhitespace=false,         
	breaklines=true,                 
	captionpos=b,                    
	keepspaces=true,                 
	%numbers=left,                    
	numbersep=5pt,                  
	showspaces=false,                
	showstringspaces=false,
	showtabs=false,                  
	tabsize=2
}
\lstset{style=myCodeStyle}
\renewcommand{\lstlistingname}{Código}%Muda o termo "Listing":  Listing -> Código
\renewcommand{\lstlistlistingname}{Lista de \lstlistingname s}% List of Listings -> Lista de Códigos


\newcommand{\unit}[1]{\ensuremath{\, \mathrm{#1}}} % comando pra tirar do math mode a unidade de medida

\def\tamanhoGrafico{0.8}

\newcommand{\ohm}{\unit{\Omega} }
\newcommand{\kohm}{\unit{k\Omega}}

\title{
	Laboratório Digital A - PCS3335
	\\ Planejamento da Experiência 2
}

\author{
	Alunos:
	\\ Rafael Hipólito Rodrigues - N.USP: 14604308
	\\ {{PERSON-2}} - N.USP: {{PERSON-ID-2}}
	\\ \\
	Turma 3 - Bancada A6
	\\ Professor: Bruno de Carvalho Albertini
}

%\date{}

\begin{document}
	
	\maketitle
	
	\section*{Execução}
	
	\begin{lstlisting}[language=vhdl, caption={VHDL do entity mainEntityExp2}, label={codigo}]
entity mainEntityExp2 is
	port (
	--input counter
	counter_clock, counter_reset, counter_enable, counter_load, counter_up : in std_logic;
	
	--input/output shiftregister
	reg_clock, reg_reset, reg_serial_i : in std_logic;
	reg_loadOrShift : in std_logic_vector( 1 downto 0 );
	reg_serial_o_r, reg_serial_o_l : out std_logic;
	
	--entrada em paralelo das chaves
	keys_data_i : in std_logic_vector( 9 downto 0 );
	
	--origem do reg_data_i
	reg_source_data_i: in std_logic; -- 0 -> reg_data_i <= chaves; 1 -> reg_data_i <= contador
	
	--output display-7-seg counter
	counter_seg2, counter_seg1, counter_seg0: out std_logic_vector(6 downto 0);
	
	--output display-7-seg shiftregister
	reg_seg2, reg_seg1, reg_seg0: out std_logic_vector(6 downto 0)
	);
end entity;
	\end{lstlisting}
	
	\subsection*{Diagramas RTL}
	
	\begin{figure}[H]
		\centering
		\includegraphics[width=0.7\linewidth]{images/mainEntityExp2RTL.png}
		\caption{Diagrama RTL do mainEntityExp2.}
		\label{fig:mainEntityExp2}
	\end{figure}
	
	\begin{figure}[H]
		\centering
		\includegraphics[width=0.7\linewidth]{images/counterRTL.png}
		\caption{Diagrama RTL do counter.}
		\label{fig:counterRTL}
	\end{figure}
	
	\begin{figure}[H]
		\centering
		\includegraphics[width=0.7\linewidth]{images/regRTL.png}
		\caption{Diagrama RTL do shiftregister.}
		\label{fig:regRTL}
	\end{figure}
	
	\subsection*{Pinagem}
	
	\subsection*{Testes}

	\begin{center}
		\begin{tabular}{ |c|c|c| }
			\hline
			Valor dos Switchs(Binário) & Valor esperado(Hexa) & Valor obtido(Hexa) \\
			\hline
			00 0000 0000 & 000000 & 000000 \\
			\hline
			00 0000 0001 & 100001 & 100001 \\
			\hline
			00 0000 0010 & 200002 & 200002 \\
			\hline
			00 0000 1111 & F0000F & F0000F \\
			\hline
			00 0001 0000 & F1001F & F1001F \\
			\hline
			00 1111 1111 & FF00FF & FF00FF \\
			\hline
			01 0000 0000 & FF11FF & FF11FF \\
			\hline
			11 1111 1111 & FF33FF & FF33FF \\
			\hline
		\end{tabular}
	\end{center}
	
	Todos os resultados obtidos foram conforme o esperado.
	
	
	
\end{document}
